% =====================================================================
\chapter{Funciones reales elementales y sus gráficas}
\label{tema:1}
% =====================================================================

Todo el cálculo trata de funciones reales de una variable real. Este tema fija el vocabulario (dominio, imagen, gráfica, monotonía, paridad, inversa), repasa el catálogo de funciones elementales que se usarán como ejemplos en el resto del módulo y presenta las funciones de activación de las redes neuronales como funciones elementales más. Se siguen las §1.1--1.5 de Stewart~\cite{stewart}, las §1.1--1.5 de OpenStax~\cite{osc1} y los capítulos~3 y~4 de Spivak~\cite{spivak}; la parte conjuntista (inyectividad, composición e inversa) se estudió en el módulo~0.2.

\section{Funciones reales de variable real}

\begin{definicion}[Función real, dominio, imagen y gráfica]\label{def:funcion}
Una \emph{función real de variable real} es una función $f\colon D\to\R$ con $D\subseteq\R$. El conjunto $D=\dom f$ es su \emph{dominio}; $f(D)=\{f(x):x\in D\}$ es su \emph{imagen} o recorrido, y su \emph{gráfica} es el conjunto $\{(x,f(x)):x\in D\}\subseteq\R^{2}$. Cuando una función se da por una fórmula sin indicar el dominio, se sobrentiende el \emph{dominio natural}: el mayor conjunto de reales en el que la fórmula tiene sentido.
\end{definicion}

\addtocounter{definicion}{1}%

Las funciones se suman, se multiplican y se dividen punto a punto, en la intersección de los dominios (y, para el cociente, donde el denominador no se anula). La \emph{composición} $(g\circ f)(x)=g(f(x))$ está definida en $\{x\in\dom f: f(x)\in\dom g\}$.

\begin{definicion}[Acotación, monotonía, paridad y periodicidad]\label{def:propiedades}
Sea $f\colon D\to\R$.
\begin{enumerate}
  \item $f$ está \emph{acotada superiormente} si existe $M$ con $f(x)\le M$ para todo $x\in D$; \emph{inferiormente}, si existe $m$ con $f(x)\ge m$; \emph{acotada}, si lo está de ambas formas, es decir, si existe $K$ con $\abs{f(x)}\le K$ para todo $x\in D$.
  \item $f$ es \emph{creciente} si $x<y\Rightarrow f(x)\le f(y)$, y \emph{estrictamente creciente} si $x<y\Rightarrow f(x)<f(y)$, para todos $x,y\in D$. Análogamente, (estrictamente) \emph{decreciente}. Se llama \emph{monótona} a una función creciente o decreciente.
  \item Si $D$ es simétrico ($x\in D\Rightarrow -x\in D$), $f$ es \emph{par} si $f(-x)=f(x)$ e \emph{impar} si $f(-x)=-f(x)$ para todo $x\in D$.
  \item $f$ es \emph{periódica} de periodo $T>0$ si $x+T\in D$ y $f(x+T)=f(x)$ para todo $x\in D$.
\end{enumerate}
\end{definicion}

La gráfica de una función par es simétrica respecto del eje de ordenadas, y la de una impar, respecto del origen. Una función impar definida en~$0$ cumple $f(0)=-f(0)$, luego $f(0)=0$.

\begin{proposicion}[Transformaciones de gráficas]\label{prop:transformaciones}
Sea $f$ una función, $c\in\R$ y $k>0$. La gráfica de
\begin{enumerate}
  \item $y=f(x)+c$ es la de $f$ trasladada $c$ unidades en vertical;
  \item $y=f(x-c)$ es la de $f$ trasladada $c$ unidades en horizontal (hacia la derecha si $c>0$);
  \item $y=kf(x)$ es la de $f$ dilatada verticalmente por el factor $k$;
  \item $y=f(kx)$ es la de $f$ comprimida horizontalmente por el factor $k$ (dilatada si $k<1$);
  \item $y=-f(x)$ e $y=f(-x)$ son sus simétricas respecto del eje de abscisas y del de ordenadas.
\end{enumerate}
\end{proposicion}

\section{El catálogo de funciones elementales}

\begin{definicion}[Funciones elementales básicas]\label{def:elementales}
Se llaman \emph{elementales} las funciones que se obtienen, mediante sumas, productos, cocientes, composiciones e inversas, a partir de:
\begin{enumerate}
  \item los \emph{polinomios} $p(x)=a_nx^{n}+\dots+a_1x+a_0$ y las \emph{funciones racionales} $p/q$ (cocientes de polinomios, en los puntos donde $q\ne0$);
  \item las \emph{potencias} $x^{a}$ ($x>0$, $a\in\R$) y las raíces $\sqrt[n]{x}$;
  \item la \emph{exponencial} $\exp(x)=e^{x}$, biyección creciente de $\R$ en $(0,\infty)$, y su inversa, el \emph{logaritmo} $\ln\colon(0,\infty)\to\R$;
  \item las funciones trigonométricas $\sin$, $\cos$, $\tan$ y sus inversas $\arcsin$, $\arccos$, $\arctan$ (módulo~0.4).
\end{enumerate}
\end{definicion}

Se recuerdan las propiedades de la exponencial y del logaritmo (módulo~0.3): $e^{x+y}=e^{x}e^{y}$, $e^{0}=1$, $\ln(xy)=\ln x+\ln y$, $\ln(x^{a})=a\ln x$, $e^{\ln x}=x$ para $x>0$ y $\ln(e^{x})=x$ para todo $x$; además $a^{x}=e^{x\ln a}$ para $a>0$. El valor absoluto $\abs{x}=\max(x,-x)$ y la parte entera $\floor{x}$ (el mayor entero $\le x$) completan el repertorio de ejemplos.

\begin{definicion}[Funciones hiperbólicas]\label{def:hiperbolicas}
Para $x\in\R$ se definen
\[
\sinh x=\frac{e^{x}-e^{-x}}{2},\qquad \cosh x=\frac{e^{x}+e^{-x}}{2},\qquad \tanh x=\frac{\sinh x}{\cosh x}=\frac{e^{x}-e^{-x}}{e^{x}+e^{-x}}.
\]
\end{definicion}

\begin{proposicion}[Propiedades de las funciones hiperbólicas]\label{prop:hiperbolicas}
Para todo $x\in\R$:
\begin{enumerate}
  \item $\cosh$ es par, $\cosh x\ge1$, y $\sinh$ y $\tanh$ son impares.
  \item $\cosh^{2}x-\sinh^{2}x=1$ para todo $x$.
  \item $\cosh x+\sinh x=e^{x}$ y $\cosh x-\sinh x=e^{-x}$.
  \item $\abs{\tanh x}<1$, y $\tanh$ es estrictamente creciente de $\R$ en $(-1,1)$.
\end{enumerate}
\end{proposicion}

El nombre procede de (ii): el punto $(\cosh t,\sinh t)$ recorre la rama derecha de la hipérbola $x^{2}-y^{2}=1$, como $(\cos t,\sin t)$ recorre la circunferencia $x^{2}+y^{2}=1$.

\section{Funciones de activación}

Una neurona artificial calcula una combinación lineal de sus entradas y le aplica una función no lineal fija, la \emph{función de activación}. Las más usadas son funciones elementales (o casi) y sirven como ejemplos recurrentes en este módulo.

\begin{definicion}[Funciones de activación]\label{def:activaciones}
Para $x\in\R$ se definen
\begin{enumerate}
  \item la \emph{sigmoide} o función logística: $\sigma(x)=\dfrac{1}{1+e^{-x}}$;
  \item la \emph{tangente hiperbólica} $\tanh x$ (\cref{def:hiperbolicas});
  \item la \emph{unidad lineal rectificada}: $\ReLU(x)=\max(0,x)$;
  \item la \emph{softplus}: $\softplus(x)=\ln(1+e^{x})$.
\end{enumerate}
\end{definicion}

\addtocounter{figure}{1}%

\begin{proposicion}[Propiedades de la sigmoide]\label{prop:sigmoide}
La función sigmoide cumple:
\begin{enumerate}
  \item $0<\sigma(x)<1$ para todo $x$, y $\sigma(0)=\tfrac12$.
  \item $\sigma(-x)=1-\sigma(x)$; en consecuencia, la gráfica es simétrica respecto del punto $(0,\tfrac12)$.
  \item $\sigma$ es estrictamente creciente y biyectiva de $\R$ en $(0,1)$, con inversa $\logit(p)=\ln\dfrac{p}{1-p}$.
  \item $\tanh x=2\sigma(2x)-1$.
\end{enumerate}
\end{proposicion}

La propiedad (iv) dice que $\tanh$ es una sigmoide reescalada: misma forma, imagen $(-1,1)$ en lugar de $(0,1)$ y pendiente cuatro veces mayor en el origen (\cref{tema:4}). En la regresión logística (módulo~8), $\sigma$ transforma una puntuación real en una probabilidad, y $\logit$ deshace esa transformación: el logit de una probabilidad es el logaritmo de la razón de probabilidades $p/(1-p)$, de donde procede el término \emph{logits} con que se designan las salidas de una red antes de la softmax.

\begin{proposicion}[ReLU y softplus]\label{prop:relu}
Para todo $x\in\R$:
\begin{enumerate}
  \item $\ReLU(x)=\dfrac{x+\abs{x}}{2}$ y $\ReLU(x)-\ReLU(-x)=x$;
  \item $\softplus(x)-\softplus(-x)=x$;
  \item $\ReLU(x)<\softplus(x)\le\ReLU(x)+\ln 2$, con igualdad a la derecha solo en $x=0$.
\end{enumerate}
\end{proposicion}

La softplus es, por tanto, una versión «suavizada» de la ReLU que nunca se separa de ella más de $\ln 2\approx0{,}693$ (\cref{fig:activaciones}). Otras activaciones modernas, como la GELU, $\GELU(x)=x\,\Phi(x)$ con $\Phi$ la función de distribución de la normal estándar, no son elementales; se estudian en el módulo~9.

\section{Función inversa}

\begin{proposicion}[Monotonía estricta e inversa]\label{prop:inversa-monotona}
Sea $f\colon D\to\R$ estrictamente creciente. Entonces $f$ es inyectiva, su inversa $f^{-1}\colon f(D)\to D$ es estrictamente creciente, y la gráfica de $f^{-1}$ es la simétrica de la de $f$ respecto de la recta $y=x$.
\end{proposicion}

\addtocounter{definicion}{1}%

\addtocounter{definicion}{1}%

\addtocounter{definicion}{1}%

